\documentclass[aspectratio=169,10pt]{beamer}
\usetheme{Madrid}
\usecolortheme{default}
\usepackage[utf8]{inputenc}
\usepackage[T1]{fontenc}
\usepackage{lmodern}
\usepackage{amsmath,amssymb}
\usepackage{enumitem}
\usepackage{xcolor}

\setbeamertemplate{navigation symbols}{}
\setbeamertemplate{footline}[frame number]
\setbeamertemplate{itemize items}[circle]
\setbeamerfont{frametitle}{size=\large,series=\bfseries}

\title{CIA Part 1 -- Ethics \& Professionalism}
\subtitle{100 Hard Questions with Detailed Explanations}
\author{Study Deck}
\date{\today}

\begin{document}
	
	\begin{frame}
		\titlepage
	\end{frame}
	
	\begin{frame}{Format \& Coverage}
		\small
		\begin{itemize}
			\item 100 hard questions on Ethics \& Professionalism.
			\item Each question is followed by a detailed explanation slide.
			\item Coverage: IIA Code of Ethics (Principles \& Rules of Conduct),
			objectivity, conflicts of interest, gifts/hospitality, confidentiality,
			competency, CPD, professional skepticism, due care, independence vs.
			objectivity, impairment disclosure, whistleblowing, ethical decision
			frameworks, and disciplinary actions.
		\end{itemize}
	\end{frame}
	
	% =========================================================
	% Q1
	% =========================================================
	\begin{frame}[t]{Q1}
		\small
		\textbf{Question:} The IIA Code of Ethics applies to:
		\begin{enumerate}[label=\Alph*.,itemsep=1pt,topsep=2pt]
			\item Only CIA-certified internal auditors.
			\item All individuals and entities that provide internal audit services.
			\item Only CAEs.
			\item Only members of the IIA.
		\end{enumerate}
		\vspace{4pt}
		\textbf{Correct Answer: B}
	\end{frame}
	
	\begin{frame}[t]{Q1 -- Explanation}
		\small
		\textbf{Correct: B.} The Code applies to all individuals and entities performing internal audit services, including IIA members, CIA holders, and those who provide internal audit services under the Standards. \\[4pt]
		\textbf{A wrong:} Non-certified internal auditors are also bound. \\[2pt]
		\textbf{C wrong:} It applies to all internal auditors, not just CAEs. \\[2pt]
		\textbf{D wrong:} It extends beyond IIA membership to anyone providing internal audit services.
	\end{frame}
	
	% =========================================================
	% Q2
	% =========================================================
	\begin{frame}[t]{Q2}
		\small
		\textbf{Question:} How many principles are in the IIA Code of Ethics?
		\begin{enumerate}[label=\Alph*.,itemsep=1pt,topsep=2pt]
			\item Two.
			\item Three.
			\item Four.
			\item Five.
		\end{enumerate}
		\vspace{4pt}
		\textbf{Correct Answer: C}
	\end{frame}
	
	\begin{frame}[t]{Q2 -- Explanation}
		\small
		\textbf{Correct: C.} Four principles: Integrity, Objectivity, Confidentiality, and Competency. \\[4pt]
		\textbf{A wrong:} Two is too few. \\[2pt]
		\textbf{B wrong:} Three is incorrect. \\[2pt]
		\textbf{D wrong:} Five is incorrect -- a common distractor.
	\end{frame}
	
	% =========================================================
	% Q3
	% =========================================================
	\begin{frame}[t]{Q3}
		\small
		\textbf{Question:} Which is NOT one of the four principles of the IIA Code of Ethics?
		\begin{enumerate}[label=\Alph*.,itemsep=1pt,topsep=2pt]
			\item Integrity.
			\item Objectivity.
			\item Independence.
			\item Confidentiality.
		\end{enumerate}
		\vspace{4pt}
		\textbf{Correct Answer: C}
	\end{frame}
	
	\begin{frame}[t]{Q3 -- Explanation}
		\small
		\textbf{Correct: C.} Independence is a Standard/attribute concept, not one of the four Code principles. \\[4pt]
		\textbf{A wrong:} Integrity is a principle. \\[2pt]
		\textbf{B wrong:} Objectivity is a principle. \\[2pt]
		\textbf{D wrong:} Confidentiality is a principle.
	\end{frame}
	
	% =========================================================
	% Q4
	% =========================================================
	\begin{frame}[t]{Q4}
		\small
		\textbf{Question:} The Code of Ethics has two essential components:
		\begin{enumerate}[label=\Alph*.,itemsep=1pt,topsep=2pt]
			\item Principles and Rules of Conduct.
			\item Laws and Regulations.
			\item Standards and Guidelines.
			\item Policies and Procedures.
		\end{enumerate}
		\vspace{4pt}
		\textbf{Correct Answer: A}
	\end{frame}
	
	\begin{frame}[t]{Q4 -- Explanation}
		\small
		\textbf{Correct: A.} The Code consists of Principles (conceptual) and Rules of Conduct (behavioral expectations). \\[4pt]
		\textbf{B wrong:} Laws/regulations are external to the Code. \\[2pt]
		\textbf{C wrong:} Standards are separate from the Code. \\[2pt]
		\textbf{D wrong:} Policies/procedures are organizational.
	\end{frame}
	
	% =========================================================
	% Q5
	% =========================================================
	\begin{frame}[t]{Q5}
		\small
		\textbf{Question:} Which principle requires internal auditors to perform work with honesty, diligence, and responsibility?
		\begin{enumerate}[label=\Alph*.,itemsep=1pt,topsep=2pt]
			\item Integrity.
			\item Objectivity.
			\item Confidentiality.
			\item Competency.
		\end{enumerate}
		\vspace{4pt}
		\textbf{Correct Answer: A}
	\end{frame}
	
	\begin{frame}[t]{Q5 -- Explanation}
		\small
		\textbf{Correct: A.} Integrity requires honesty, diligence, and responsibility; observing the law; and not knowingly being part of illegal or discreditable activity. \\[4pt]
		\textbf{B wrong:} Objectivity concerns bias and conflicts. \\[2pt]
		\textbf{C wrong:} Confidentiality concerns information protection. \\[2pt]
		\textbf{D wrong:} Competency concerns skills and CPD.
	\end{frame}
	
	% =========================================================
	% Q6
	% =========================================================
	\begin{frame}[t]{Q6}
		\small
		\textbf{Question:} An internal auditor must not accept anything that impairs or appears to impair professional judgment. This relates to:
		\begin{enumerate}[label=\Alph*.,itemsep=1pt,topsep=2pt]
			\item Integrity.
			\item Objectivity.
			\item Confidentiality.
			\item Competency.
		\end{enumerate}
		\vspace{4pt}
		\textbf{Correct Answer: B}
	\end{frame}
	
	\begin{frame}[t]{Q6 -- Explanation}
		\small
		\textbf{Correct: B.} Objectivity prohibits accepting gifts, participating in activities, or having relationships that impair -- or appear to impair -- unbiased judgment. \\[4pt]
		\textbf{A wrong:} Integrity is about honesty, not gifts. \\[2pt]
		\textbf{C wrong:} Confidentiality concerns information. \\[2pt]
		\textbf{D wrong:} Competency concerns skills.
	\end{frame}
	
	% =========================================================
	% Q7
	% =========================================================
	\begin{frame}[t]{Q7}
		\small
		\textbf{Question:} An auditor must be prudent in the use and protection of information acquired. This relates to:
		\begin{enumerate}[label=\Alph*.,itemsep=1pt,topsep=2pt]
			\item Integrity.
			\item Objectivity.
			\item Confidentiality.
			\item Competency.
		\end{enumerate}
		\vspace{4pt}
		\textbf{Correct Answer: C}
	\end{frame}
	
	\begin{frame}[t]{Q7 -- Explanation}
		\small
		\textbf{Correct: C.} Confidentiality requires prudence in the use/protection of information and prohibits disclosing information without authority (unless legally/ professionally required). \\[4pt]
		\textbf{A wrong:} Integrity concerns honesty. \\[2pt]
		\textbf{B wrong:} Objectivity concerns bias. \\[2pt]
		\textbf{D wrong:} Competency concerns skills.
	\end{frame}
	
	% =========================================================
	% Q8
	% =========================================================
	\begin{frame}[t]{Q8}
		\small
		\textbf{Question:} Auditors must engage only in services for which they have the necessary knowledge, skills, and experience. This relates to:
		\begin{enumerate}[label=\Alph*.,itemsep=1pt,topsep=2pt]
			\item Integrity.
			\item Objectivity.
			\item Confidentiality.
			\item Competency.
		\end{enumerate}
		\vspace{4pt}
		\textbf{Correct Answer: D}
	\end{frame}
	
	\begin{frame}[t]{Q8 -- Explanation}
		\small
		\textbf{Correct: D.} Competency requires the knowledge, skills, and experience needed, and performing services only when competent (or obtaining competent advice). \\[4pt]
		\textbf{A wrong:} Integrity concerns honesty. \\[2pt]
		\textbf{B wrong:} Objectivity concerns bias. \\[2pt]
		\textbf{C wrong:} Confidentiality concerns information.
	\end{frame}
	
	% =========================================================
	% Q9
	% =========================================================
	\begin{frame}[t]{Q9}
		\small
		\textbf{Question:} An internal auditor is offered a modest gift by an audit client that could be seen as influencing judgment. The auditor should:
		\begin{enumerate}[label=\Alph*.,itemsep=1pt,topsep=2pt]
			\item Accept it, since it is modest.
			\item Decline it, or disclose to the CAE.
			\item Accept it but disclose after the engagement.
			\item Give it to a colleague.
		\end{enumerate}
		\vspace{4pt}
		\textbf{Correct Answer: B}
	\end{frame}
	
	\begin{frame}[t]{Q9 -- Explanation}
		\small
		\textbf{Correct: B.} The auditor should decline or disclose to the CAE. Even the appearance of impairment violates Objectivity. \\[4pt]
		\textbf{A wrong:} Modesty does not remove the appearance of impairment. \\[2pt]
		\textbf{C wrong:} Disclosure must be before acceptance, not after. \\[2pt]
		\textbf{D wrong:} Redirecting does not solve the ethical issue.
	\end{frame}
	
	% =========================================================
	% Q10
	% =========================================================
	\begin{frame}[t]{Q10}
		\small
		\textbf{Question:} Which of the following is a Rules of Conduct requirement under Objectivity?
		\begin{enumerate}[label=\Alph*.,itemsep=1pt,topsep=2pt]
			\item Shall not participate in any activity that may impair unbiased assessment.
			\item Shall be prudent in the use of information.
			\item Shall engage only in services for which they are competent.
			\item Shall perform work with honesty.
		\end{enumerate}
		\vspace{4pt}
		\textbf{Correct Answer: A}
	\end{frame}
	
	\begin{frame}[t]{Q10 -- Explanation}
		\small
		\textbf{Correct: A.} Objectivity's rules prohibit activities/relationships that impair unbiased assessment and require disclosure of material facts. \\[4pt]
		\textbf{B wrong:} That's Confidentiality. \\[2pt]
		\textbf{C wrong:} That's Competency. \\[2pt]
		\textbf{D wrong:} That's Integrity.
	\end{frame}
	
	% =========================================================
	% Q11
	% =========================================================
	\begin{frame}[t]{Q11}
		\small
		\textbf{Question:} An auditor discovers fraud during an engagement. Under Integrity, they must:
		\begin{enumerate}[label=\Alph*.,itemsep=1pt,topsep=2pt]
			\item Ignore it if immaterial.
			\item Report it appropriately and not be part of illegal activity.
			\item Hide it to protect the company.
			\item Discuss it only with the suspect.
		\end{enumerate}
		\vspace{4pt}
		\textbf{Correct Answer: B}
	\end{frame}
	
	\begin{frame}[t]{Q11 -- Explanation}
		\small
		\textbf{Correct: B.} Integrity requires reporting and refusing to be part of illegal/discreditable activity. \\[4pt]
		\textbf{A wrong:} Fraud is never ignored. \\[2pt]
		\textbf{C wrong:} Hiding is unethical. \\[2pt]
		\textbf{D wrong:} Confronting the suspect can compromise the investigation.
	\end{frame}
	
	% =========================================================
	% Q12
	% =========================================================
	\begin{frame}[t]{Q12}
		\small
		\textbf{Question:} Which is the correct order of the ethical decision-making process?
		\begin{enumerate}[label=\Alph*.,itemsep=1pt,topsep=2pt]
			\item Identify facts, define ethical issue, identify principles, evaluate options, act.
			\item Act first, then justify.
			\item Ignore the issue if minor.
			\item Ask a colleague and act on their advice.
		\end{enumerate}
		\vspace{4pt}
		\textbf{Correct Answer: A}
	\end{frame}
	
	\begin{frame}[t]{Q12 -- Explanation}
		\small
		\textbf{Correct: A.} A structured ethical framework: gather facts, define the issue, apply the Code's principles, evaluate alternatives, decide, and act. \\[4pt]
		\textbf{B wrong:} Acting first is reckless. \\[2pt]
		\textbf{C wrong:} Ignoring is not a framework. \\[2pt]
		\textbf{D wrong:} Colleague input helps but isn't the framework.
	\end{frame}
	
	% =========================================================
	% Q13
	% =========================================================
	\begin{frame}[t]{Q13}
		\small
		\textbf{Question:} The CAE is asked to approve an audit of a process she previously managed. She should:
		\begin{enumerate}[label=\Alph*.,itemsep=1pt,topsep=2pt]
			\item Proceed, since she knows the process.
			\item Decline to lead; assign to another auditor and disclose.
			\item Delegate to a subordinate without disclosure.
			\item Ignore the conflict.
		\end{enumerate}
		\vspace{4pt}
		\textbf{Correct Answer: B}
	\end{frame}
	
	\begin{frame}[t]{Q13 -- Explanation}
		\small
		\textbf{Correct: B.} Objectivity is impaired; the CAE must disclose and assign the engagement to an auditor without the conflict. \\[4pt]
		\textbf{A wrong:} Prior involvement impairs objectivity. \\[2pt]
		\textbf{C wrong:} Delegation without disclosure hides the conflict. \\[2pt]
		\textbf{D wrong:} Conflicts cannot be ignored.
	\end{frame}
	
	% =========================================================
	% Q14
	% =========================================================
	\begin{frame}[t]{Q14}
		\small
		\textbf{Question:} Which best describes a conflict of interest?
		\begin{enumerate}[label=\Alph*.,itemsep=1pt,topsep=2pt]
			\item A situation where an auditor's personal interests could improperly influence judgment.
			\item A disagreement with management.
			\item A difference of opinion on audit scope.
			\item A budget dispute.
		\end{enumerate}
		\vspace{4pt}
		\textbf{Correct Answer: A}
	\end{frame}
	
	\begin{frame}[t]{Q14 -- Explanation}
		\small
		\textbf{Correct: A.} A conflict of interest exists when personal, financial, or other interests could improperly influence professional judgment. \\[4pt]
		\textbf{B wrong:} Disagreement is not automatically a conflict. \\[2pt]
		\textbf{C wrong:} Scope differences are professional, not conflicts. \\[2pt]
		\textbf{D wrong:} Budget disputes are administrative.
	\end{frame}
	
	% =========================================================
	% Q15
	% =========================================================
	\begin{frame}[t]{Q15}
		\small
		\textbf{Question:} Which is the best way to manage a conflict of interest?
		\begin{enumerate}[label=\Alph*.,itemsep=1pt,topsep=2pt]
			\item Ignore it if minor.
			\item Disclose it and recuse if necessary.
			\item Hide it.
			\item Delegate without disclosure.
		\end{enumerate}
		\vspace{4pt}
		\textbf{Correct Answer: B}
	\end{frame}
	
	\begin{frame}[t]{Q15 -- Explanation}
		\small
		\textbf{Correct: B.} Disclosure and, if needed, recusal are the standard responses. \\[4pt]
		\textbf{A wrong:} Ignoring is not acceptable. \\[2pt]
		\textbf{C wrong:} Hiding violates the Code. \\[2pt]
		\textbf{D wrong:} Delegation without disclosure hides the issue.
	\end{frame}
	
	% =========================================================
	% Q16
	% =========================================================
	\begin{frame}[t]{Q16}
		\small
		\textbf{Question:} An auditor's spouse owns shares in a company the auditor is assigned to audit. The auditor should:
		\begin{enumerate}[label=\Alph*.,itemsep=1pt,topsep=2pt]
			\item Proceed, since the spouse owns the shares.
			\item Disclose and decline the assignment.
			\item Sell the shares after the audit.
			\item Ignore it.
		\end{enumerate}
		\vspace{4pt}
		\textbf{Correct Answer: B}
	\end{frame}
	
	\begin{frame}[t]{Q16 -- Explanation}
		\small
		\textbf{Correct: B.} A close family member's financial interest creates a conflict; disclose and recuse. \\[4pt]
		\textbf{A wrong:} The conflict exists regardless of who owns the shares. \\[2pt]
		\textbf{C wrong:} Selling after the fact doesn't fix the impairment during the audit. \\[2pt]
		\textbf{D wrong:} Cannot be ignored.
	\end{frame}
	
	% =========================================================
	% Q17
	% =========================================================
	\begin{frame}[t]{Q17}
		\small
		\textbf{Question:} An internal auditor was previously the manager of the department being audited. The auditor should:
		\begin{enumerate}[label=\Alph*.,itemsep=1pt,topsep=2pt]
			\item Audit the department since they know it best.
			\item Decline the assignment and disclose to the CAE.
			\item Audit only the parts they didn't manage.
			\item Review their own prior work.
		\end{enumerate}
		\vspace{4pt}
		\textbf{Correct Answer: B}
	\end{frame}
	
	\begin{frame}[t]{Q17 -- Explanation}
		\small
		\textbf{Correct: B.} Auditing one's own prior work impairs objectivity; decline and disclose. \\[4pt]
		\textbf{A wrong:} Familiarity creates bias. \\[2pt]
		\textbf{C wrong:} Partial audits still create impairment. \\[2pt]
		\textbf{D wrong:} Reviewing own work is a self-review threat.
	\end{frame}
	
	% =========================================================
	% Q18
	% =========================================================
	\begin{frame}[t]{Q18}
		\small
		\textbf{Question:} Which is the best description of "professional skepticism"?
		\begin{enumerate}[label=\Alph*.,itemsep=1pt,topsep=2pt]
			\item Assuming management is dishonest.
			\item A questioning mind and critical assessment of evidence.
			\item Ignoring evidence.
			\item Relying on prior-year conclusions.
		\end{enumerate}
		\vspace{4pt}
		\textbf{Correct Answer: B}
	\end{frame}
	
	\begin{frame}[t]{Q18 -- Explanation}
		\small
		\textbf{Correct: B.} Professional skepticism is a questioning mind, alertness to conditions suggesting misstatement, and critical evaluation of evidence. \\[4pt]
		\textbf{A wrong:} It is not cynicism. \\[2pt]
		\textbf{C wrong:} Ignoring evidence is the opposite. \\[2pt]
		\textbf{D wrong:} Prior conclusions must be re-validated.
	\end{frame}
	
	% =========================================================
	% Q19
	% =========================================================
	\begin{frame}[t]{Q19}
		\small
		\textbf{Question:} An auditor is asked by the CEO to change a finding in the report. The auditor should:
		\begin{enumerate}[label=\Alph*.,itemsep=1pt,topsep=2pt]
			\item Change it as instructed.
			\item Decline if the change would misrepresent facts, and escalate if needed.
			\item Ignore the request and issue the report.
			\item Resign.
		\end{enumerate}
		\vspace{4pt}
		\textbf{Correct Answer: B}
	\end{frame}
	
	\begin{frame}[t]{Q19 -- Explanation}
		\small
		\textbf{Correct: B.} Integrity requires refusing to misrepresent facts; escalate through the CAE/audit committee if pressured. \\[4pt]
		\textbf{A wrong:} Changing facts violates Integrity. \\[2pt]
		\textbf{C wrong:} Escalation is needed for a serious request. \\[2pt]
		\textbf{D wrong:} Resignation is a last resort.
	\end{frame}
	
	% =========================================================
	% Q20
	% =========================================================
	\begin{frame}[t]{Q20}
		\small
		\textbf{Question:} Which of the following is a violation of Confidentiality?
		\begin{enumerate}[label=\Alph*.,itemsep=1pt,topsep=2pt]
			\item Discussing audit findings in a public cafeteria.
			\item Sharing findings with the audit committee.
			\item Sharing findings with the external auditor with authorization.
			\item Disclosing information required by law.
		\end{enumerate}
		\vspace{4pt}
		\textbf{Correct Answer: A}
	\end{frame}
	
	\begin{frame}[t]{Q20 -- Explanation}
		\small
		\textbf{Correct: A.} Public discussion of confidential findings violates Confidentiality. \\[4pt]
		\textbf{B wrong:} Audit committee reporting is authorized. \\[2pt]
		\textbf{C wrong:} Authorized sharing is permitted. \\[2pt]
		\textbf{D wrong:} Legal disclosure is required.
	\end{frame}
	
	% =========================================================
	% Q21
	% =========================================================
	\begin{frame}[t]{Q21}
		\small
		\textbf{Question:} An auditor leaves a laptop with client data on a train. This most directly violates:
		\begin{enumerate}[label=\Alph*.,itemsep=1pt,topsep=2pt]
			\item Integrity.
			\item Objectivity.
			\item Confidentiality.
			\item Competency.
		\end{enumerate}
		\vspace{4pt}
		\textbf{Correct Answer: C}
	\end{frame}
	
	\begin{frame}[t]{Q21 -- Explanation}
		\small
		\textbf{Correct: C.} Failing to protect data violates the Confidentiality principle's requirement for prudence. \\[4pt]
		\textbf{A wrong:} Integrity is about honesty. \\[2pt]
		\textbf{B wrong:} Objectivity is about bias. \\[2pt]
		\textbf{D wrong:} Competency is about skills.
	\end{frame}
	
	% =========================================================
	% Q22
	% =========================================================
	\begin{frame}[t]{Q22}
		\small
		\textbf{Question:} Which best describes the IIA's position on using confidential information for personal benefit?
		\begin{enumerate}[label=\Alph*.,itemsep=1pt,topsep=2pt]
			\item Permitted if not harmful.
			\item Prohibited.
			\item Permitted if disclosed.
			\item Permitted for family.
		\end{enumerate}
		\vspace{4pt}
		\textbf{Correct Answer: B}
	\end{frame}
	
	\begin{frame}[t]{Q22 -- Explanation}
		\small
		\textbf{Correct: B.} Using confidential information for personal benefit is prohibited under Confidentiality. \\[4pt]
		\textbf{A wrong:} No personal use is permitted. \\[2pt]
		\textbf{C wrong:} Disclosure doesn't make it permissible. \\[2pt]
		\textbf{D wrong:} Family use is also prohibited.
	\end{frame}
	
	% =========================================================
	% Q23
	% =========================================================
	\begin{frame}[t]{Q23}
		\small
		\textbf{Question:} The IIA Code of Ethics requires internal auditors to report to the board any illegal activity they:
		\begin{enumerate}[label=\Alph*.,itemsep=1pt,topsep=2pt]
			\item Suspect.
			\item Are asked to participate in.
			\item Knowingly become involved in.
			\item Hear rumors about.
		\end{enumerate}
		\vspace{4pt}
		\textbf{Correct Answer: C}
	\end{frame}
	
	\begin{frame}[t]{Q23 -- Explanation}
		\small
		\textbf{Correct: C.} Integrity requires auditors not to knowingly be involved in illegal or discreditable activity; they must report it. \\[4pt]
		\textbf{A wrong:} Suspicion alone may not trigger the rule; knowledge does. \\[2pt]
		\textbf{B wrong:} Being asked is not the same as being involved. \\[2pt]
		\textbf{D wrong:} Rumors are not knowledge.
	\end{frame}
	
	% =========================================================
	% Q24
	% =========================================================
	\begin{frame}[t]{Q24}
		\small
		\textbf{Question:} An auditor is asked to perform a forensic IT audit but has no IT forensic training. Under Competency, the auditor should:
		\begin{enumerate}[label=\Alph*.,itemsep=1pt,topsep=2pt]
			\item Proceed anyway.
			\item Decline or obtain competent assistance.
			\item Ask management to do it.
			\item Skip the risky areas.
		\end{enumerate}
		\vspace{4pt}
		\textbf{Correct Answer: B}
	\end{frame}
	
	\begin{frame}[t]{Q24 -- Explanation}
		\small
		\textbf{Correct: B.} Competency requires declining or obtaining competent advice/assistance. \\[4pt]
		\textbf{A wrong:} Proceeding without competence violates the Code. \\[2pt]
		\textbf{C wrong:} Management cannot replace the auditor's responsibility. \\[2pt]
		\textbf{D wrong:} Skipping areas defeats the engagement.
	\end{frame}
	
	% =========================================================
	% Q25
	% =========================================================
	\begin{frame}[t]{Q25}
		\small
		\textbf{Question:} Which is NOT a Rule of Conduct under Competency?
		\begin{enumerate}[label=\Alph*.,itemsep=1pt,topsep=2pt]
			\item Engage only in services for which you have the necessary knowledge.
			\item Perform services in accordance with the Standards.
			\item Continually improve proficiency.
			\item Guarantee that all fraud is detected.
		\end{enumerate}
		\vspace{4pt}
		\textbf{Correct Answer: D}
	\end{frame}
	
	\begin{frame}[t]{Q25 -- Explanation}
		\small
		\textbf{Correct: D.} No auditor can guarantee fraud detection; Competency requires knowledge, Standards compliance, and CPD. \\[4pt]
		\textbf{A wrong:} That is a Competency rule. \\[2pt]
		\textbf{B wrong:} That is a Competency rule. \\[2pt]
		\textbf{C wrong:} That is a Competency rule.
	\end{frame}
	
	% =========================================================
	% Q26
	% =========================================================
	\begin{frame}[t]{Q26}
		\small
		\textbf{Question:} Which best describes the relationship between Independence and Objectivity?
		\begin{enumerate}[label=\Alph*.,itemsep=1pt,topsep=2pt]
			\item They are the same thing.
			\item Independence is an organizational attribute; objectivity is an individual state of mind.
			\item Objectivity is organizational; independence is individual.
			\item Neither matters for internal audit.
		\end{enumerate}
		\vspace{4pt}
		\textbf{Correct Answer: B}
	\end{frame}
	
	\begin{frame}[t]{Q26 -- Explanation}
		\small
		\textbf{Correct: B.} Independence is an organizational attribute (reporting lines, freedom from interference); objectivity is an individual's unbiased mental attitude. \\[4pt]
		\textbf{A wrong:} They are distinct. \\[2pt]
		\textbf{C wrong:} The definitions are reversed. \\[2pt]
		\textbf{D wrong:} Both are essential.
	\end{frame}
	
	% =========================================================
	% Q27
	% =========================================================
	\begin{frame}[t]{Q27}
		\small
		\textbf{Question:} An auditor has a close friend in the department being audited. The auditor should:
		\begin{enumerate}[label=\Alph*.,itemsep=1pt,topsep=2pt]
			\item Proceed; friendship doesn't matter.
			\item Disclose the relationship and consider reassignment.
			\item Hide the friendship.
			\item Audit only the friend's area.
		\end{enumerate}
		\vspace{4pt}
		\textbf{Correct Answer: B}
	\end{frame}
	
	\begin{frame}[t]{Q27 -- Explanation}
		\small
		\textbf{Correct: B.} Close personal relationships can impair objectivity; disclose and consider reassignment. \\[4pt]
		\textbf{A wrong:} Friendship can bias judgment. \\[2pt]
		\textbf{C wrong:} Hiding violates the Code. \\[2pt]
		\textbf{D wrong:} Auditing the friend's area is the worst choice.
	\end{frame}
	
	% =========================================================
	% Q28
	% =========================================================
	\begin{frame}[t]{Q28}
		\small
		\textbf{Question:} Which is the best description of "due professional care"?
		\begin{enumerate}[label=\Alph*.,itemsep=1pt,topsep=2pt]
			\item Guaranteeing no errors.
			\item Exercising the care expected of a reasonably prudent auditor.
			\item Performing all possible procedures.
			\item Avoiding all risk.
		\end{enumerate}
		\vspace{4pt}
		\textbf{Correct Answer: B}
	\end{frame}
	
	\begin{frame}[t]{Q28 -- Explanation}
		\small
		\textbf{Correct: B.} Due care is the care and skill expected of a reasonably prudent and competent auditor. \\[4pt]
		\textbf{A wrong:} No guarantee is possible. \\[2pt]
		\textbf{C wrong:} Not all procedures are feasible. \\[2pt]
		\textbf{D wrong:} Risk cannot be eliminated.
	\end{frame}
	
	% =========================================================
	% Q29
	% =========================================================
	\begin{frame}[t]{Q29}
		\small
		\textbf{Question:} An auditor is pressured by a manager to omit a finding. The auditor should:
		\begin{enumerate}[label=\Alph*.,itemsep=1pt,topsep=2pt]
			\item Omit it to avoid conflict.
			\item Keep the finding and escalate if necessary.
			\item Change the finding to a suggestion.
			\item Delay the report indefinitely.
		\end{enumerate}
		\vspace{4pt}
		\textbf{Correct Answer: B}
	\end{frame}
	
	\begin{frame}[t]{Q29 -- Explanation}
		\small
		\textbf{Correct: B.} Integrity and Objectivity require keeping the finding and escalating through proper channels. \\[4pt]
		\textbf{A wrong:} Omitting violates the Code. \\[2pt]
		\textbf{C wrong:} Downgrading misrepresents the finding. \\[2pt]
		\textbf{D wrong:} Delaying is not a solution.
	\end{frame}
	
	% =========================================================
	% Q30
	% =========================================================
	\begin{frame}[t]{Q30}
		\small
		\textbf{Question:} Which is the best description of "whistleblowing" in the context of internal audit?
		\begin{enumerate}[label=\Alph*.,itemsep=1pt,topsep=2pt]
			\item Reporting illegal or unethical activity through proper channels.
			\item Leaking to the media.
			\item Ignoring wrongdoing.
			\item Confronting the suspect.
		\end{enumerate}
		\vspace{4pt}
		\textbf{Correct Answer: A}
	\end{frame}
	
	\begin{frame}[t]{Q30 -- Explanation}
		\small
		\textbf{Correct: A.} Whistleblowing is reporting illegal/unethical activity through internal channels (or regulators if internal channels fail). \\[4pt]
		\textbf{B wrong:} Media leaks are a last resort, not the standard. \\[2pt]
		\textbf{C wrong:} Ignoring is unethical. \\[2pt]
		\textbf{D wrong:} Confronting the suspect can compromise the investigation.
	\end{frame}
	
	% =========================================================
	% Q31
	% =========================================================
	\begin{frame}[t]{Q31}
		\small
		\textbf{Question:} Which is a key responsibility of the CAE regarding ethics?
		\begin{enumerate}[label=\Alph*.,itemsep=1pt,topsep=2pt]
			\item Ensuring auditors comply with the Code of Ethics.
			\item Writing the Code of Ethics.
			\item Enforcing criminal law.
			\item Setting tax policy.
		\end{enumerate}
		\vspace{4pt}
		\textbf{Correct Answer: A}
	\end{frame}
	
	\begin{frame}[t]{Q31 -- Explanation}
		\small
		\textbf{Correct: A.} The CAE must ensure the internal audit activity complies with the Code and Standards. \\[4pt]
		\textbf{B wrong:} The IIA writes the Code. \\[2pt]
		\textbf{C wrong:} Criminal enforcement is the government's job. \\[2pt]
		\textbf{D wrong:} Tax policy is unrelated.
	\end{frame}
	
	% =========================================================
	% Q32
	% =========================================================
	\begin{frame}[t]{Q32}
		\small
		\textbf{Question:} An auditor is offered a fully paid trip to a conference by a vendor the auditor may audit. The auditor should:
		\begin{enumerate}[label=\Alph*.,itemsep=1pt,topsep=2pt]
			\item Accept it.
			\item Decline or disclose to the CAE.
			\item Accept but not attend.
			\item Accept only if the vendor is not audited this year.
		\end{enumerate}
		\vspace{4pt}
		\textbf{Correct Answer: B}
	\end{frame}
	
	\begin{frame}[t]{Q32 -- Explanation}
		\small
		\textbf{Correct: B.} The offer creates an appearance of impairment; decline or disclose. \\[4pt]
		\textbf{A wrong:} Acceptance impairs objectivity. \\[2pt]
		\textbf{C wrong:} Accepting but not attending is still improper. \\[2pt]
		\textbf{D wrong:} Future audits create the conflict.
	\end{frame}
	
	% =========================================================
	% Q33
	% =========================================================
	\begin{frame}[t]{Q33}
		\small
		\textbf{Question:} Which best describes the "appearance" standard in Objectivity?
		\begin{enumerate}[label=\Alph*.,itemsep=1pt,topsep=2pt]
			\item Only actual impairments matter.
			\item Both actual and apparent impairments matter.
			\item Appearance is irrelevant.
			\item Only financial interests matter.
		\end{enumerate}
		\vspace{4pt}
		\textbf{Correct Answer: B}
	\end{frame}
	
	\begin{frame}[t]{Q33 -- Explanation}
		\small
		\textbf{Correct: B.} The Code prohibits both actual and apparent impairments of objectivity. \\[4pt]
		\textbf{A wrong:} Appearance also matters. \\[2pt]
		\textbf{C wrong:} Appearance is relevant. \\[2pt]
		\textbf{D wrong:} Non-financial interests also matter.
	\end{frame}
	
	% =========================================================
	% Q34
	% =========================================================
	\begin{frame}[t]{Q34}
		\small
		\textbf{Question:} An internal auditor is asked to serve on a committee selecting a new ERP system. Under the Code, the auditor should:
		\begin{enumerate}[label=\Alph*.,itemsep=1pt,topsep=2pt]
			\item Accept, since it improves business knowledge.
			\item Decline, as it impairs objectivity for future audits.
			\item Accept only if approved by the CEO.
			\item Accept and delegate the audit to a subordinate.
		\end{enumerate}
		\vspace{4pt}
		\textbf{Correct Answer: B}
	\end{frame}
	
	\begin{frame}[t]{Q34 -- Explanation}
		\small
		\textbf{Correct: B.} Serving on a selection committee creates a self-review threat for future audits; decline or disclose. \\[4pt]
		\textbf{A wrong:} Knowledge gains don't override objectivity. \\[2pt]
		\textbf{C wrong:} CEO approval doesn't remove the impairment. \\[2pt]
		\textbf{D wrong:} Delegation without disclosure hides the conflict.
	\end{frame}
	
	% =========================================================
	% Q35
	% =========================================================
	\begin{frame}[t]{Q35}
		\small
		\textbf{Question:} Which is the best description of "integrity" in the Code?
		\begin{enumerate}[label=\Alph*.,itemsep=1pt,topsep=2pt]
			\item Honesty, diligence, and responsibility.
			\item Bias-free judgment.
			\item Information protection.
			\item Skill maintenance.
		\end{enumerate}
		\vspace{4pt}
		\textbf{Correct Answer: A}
	\end{frame}
	
	\begin{frame}[t]{Q35 -- Explanation}
		\small
		\textbf{Correct: A.} Integrity = honesty, diligence, responsibility; observing the law; no knowing involvement in illegal activity. \\[4pt]
		\textbf{B wrong:} That's Objectivity. \\[2pt]
		\textbf{C wrong:} That's Confidentiality. \\[2pt]
		\textbf{D wrong:} That's Competency.
	\end{frame}
	
	% =========================================================
	% Q36
	% =========================================================
	\begin{frame}[t]{Q36}
		\small
		\textbf{Question:} The Code of Ethics requires internal auditors to disclose:
		\begin{enumerate}[label=\Alph*.,itemsep=1pt,topsep=2pt]
			\item All personal matters.
			\item Material facts that, if not disclosed, could distort reported information.
			\item Only financial interests.
			\item Only family relationships.
		\end{enumerate}
		\vspace{4pt}
		\textbf{Correct Answer: B}
	\end{frame}
	
	\begin{frame}[t]{Q36 -- Explanation}
		\small
		\textbf{Correct: B.} Objectivity requires disclosing material facts that could distort reporting. \\[4pt]
		\textbf{A wrong:} Not all personal matters. \\[2pt]
		\textbf{C wrong:} Non-financial facts also matter. \\[2pt]
		\textbf{D wrong:} Non-family relationships also matter.
	\end{frame}
	
	% =========================================================
	% Q37
	% =========================================================
	\begin{frame}[t]{Q37}
		\small
		\textbf{Question:} Which is the best description of "CPE" for internal auditors?
		\begin{enumerate}[label=\Alph*.,itemsep=1pt,topsep=2pt]
			\item Continuing professional education to maintain competence.
			\item A one-time certification exam.
			\item A substitute for experience.
			\item Optional for CIA holders.
		\end{enumerate}
		\vspace{4pt}
		\textbf{Correct Answer: A}
	\end{frame}
	
	\begin{frame}[t]{Q37 -- Explanation}
		\small
		\textbf{Correct: A.} CPE maintains and enhances competence, required by the Code and Standards. \\[4pt]
		\textbf{B wrong:} The exam is initial, not ongoing. \\[2pt]
		\textbf{C wrong:} CPE complements, not replaces, experience. \\[2pt]
		\textbf{D wrong:} CPE is required for CIA holders.
	\end{frame}
	
	% =========================================================
	% Q38
	% =========================================================
	\begin{frame}[t]{Q38}
		\small
		\textbf{Question:} Which best describes the IIA's position on receiving gifts from clients?
		\begin{enumerate}[label=\Alph*.,itemsep=1pt,topsep=2pt]
			\item Always prohibited.
			\item Prohibited if they impair or appear to impair objectivity.
			\item Always permitted.
			\item Permitted if small.
		\end{enumerate}
		\vspace{4pt}
		\textbf{Correct Answer: B}
	\end{frame}
	
	\begin{frame}[t]{Q38 -- Explanation}
		\small
		\textbf{Correct: B.} Gifts are prohibited when they impair or appear to impair objectivity; some organizations set a de minimis threshold. \\[4pt]
		\textbf{A wrong:} Not always prohibited. \\[2pt]
		\textbf{C wrong:} Not always permitted. \\[2pt]
		\textbf{D wrong:} Even small gifts can create appearance issues.
	\end{frame}
	
	% =========================================================
	% Q39
	% =========================================================
	\begin{frame}[t]{Q39}
		\small
		\textbf{Question:} An auditor is asked to audit a process she designed last year. She should:
		\begin{enumerate}[label=\Alph*.,itemsep=1pt,topsep=2pt]
			\item Proceed; one year is enough time.
			\item Decline and disclose the self-review threat.
			\item Audit only the parts she didn't design.
			\item Ask a colleague to review her work.
		\end{enumerate}
		\vspace{4pt}
		\textbf{Correct Answer: B}
	\end{frame}
	
	\begin{frame}[t]{Q39 -- Explanation}
		\small
		\textbf{Correct: B.} Designing a process creates a self-review threat; decline and disclose. \\[4pt]
		\textbf{A wrong:} Time doesn't remove the impairment. \\[2pt]
		\textbf{C wrong:} Partial audits still create impairment. \\[2pt]
		\textbf{D wrong:} Colleague review doesn't cure the conflict.
	\end{frame}
	
	% =========================================================
	% Q40
	% =========================================================
	\begin{frame}[t]{Q40}
		\small
		\textbf{Question:} Which best describes "objectivity" in the Code?
		\begin{enumerate}[label=\Alph*.,itemsep=1pt,topsep=2pt]
			\item Honesty.
			\item Unbiased mental attitude.
			\item Information protection.
			\item Skill maintenance.
		\end{enumerate}
		\vspace{4pt}
		\textbf{Correct Answer: B}
	\end{frame}
	
	\begin{frame}[t]{Q40 -- Explanation}
		\small
		\textbf{Correct: B.} Objectivity is an unbiased mental attitude -- no conflicts, no gifts that impair judgment, no undisclosed material facts. \\[4pt]
		\textbf{A wrong:} That's Integrity. \\[2pt]
		\textbf{C wrong:} That's Confidentiality. \\[2pt]
		\textbf{D wrong:} That's Competency.
	\end{frame}
	
	% =========================================================
	% Q41
	% =========================================================
	\begin{frame}[t]{Q41}
		\small
		\textbf{Question:} Which best describes "confidentiality" in the Code?
		\begin{enumerate}[label=\Alph*.,itemsep=1pt,topsep=2pt]
			\item Honesty.
			\item Unbiased judgment.
			\item Prudent use and protection of information.
			\item Skill maintenance.
		\end{enumerate}
		\vspace{4pt}
		\textbf{Correct Answer: C}
	\end{frame}
	
	\begin{frame}[t]{Q41 -- Explanation}
		\small
		\textbf{Correct: C.} Confidentiality = prudent use/protection of information; no unauthorized disclosure; no personal benefit. \\[4pt]
		\textbf{A wrong:} That's Integrity. \\[2pt]
		\textbf{B wrong:} That's Objectivity. \\[2pt]
		\textbf{D wrong:} That's Competency.
	\end{frame}
	
	% =========================================================
	% Q42
	% =========================================================
	\begin{frame}[t]{Q42}
		\small
		\textbf{Question:} Which best describes "competency" in the Code?
		\begin{enumerate}[label=\Alph*.,itemsep=1pt,topsep=2pt]
			\item Honesty.
			\item Unbiased judgment.
			\item Information protection.
			\item Knowledge, skills, and experience; CPD.
		\end{enumerate}
		\vspace{4pt}
		\textbf{Correct Answer: D}
	\end{frame}
	
	\begin{frame}[t]{Q42 -- Explanation}
		\small
		\textbf{Correct: D.} Competency = knowledge, skills, experience; engage only in services you can perform; CPD. \\[4pt]
		\textbf{A wrong:} That's Integrity. \\[2pt]
		\textbf{B wrong:} That's Objectivity. \\[2pt]
		\textbf{C wrong:} That's Confidentiality.
	\end{frame}
	
	% =========================================================
	% Q43
	% =========================================================
	\begin{frame}[t]{Q43}
		\small
		\textbf{Question:} An internal auditor discovers that a colleague has violated the Code of Ethics. The auditor should:
		\begin{enumerate}[label=\Alph*.,itemsep=1pt,topsep=2pt]
			\item Ignore it.
			\item Report it through appropriate channels.
			\item Confront the colleague publicly.
			\item Cover for the colleague.
		\end{enumerate}
		\vspace{4pt}
		\textbf{Correct Answer: B}
	\end{frame}
	
	\begin{frame}[t]{Q43 -- Explanation}
		\small
		\textbf{Correct: B.} Report ethical violations through proper channels (CAE, audit committee, IIA). \\[4pt]
		\textbf{A wrong:} Ignoring violates Integrity. \\[2pt]
		\textbf{C wrong:} Public confrontation is unprofessional. \\[2pt]
		\textbf{D wrong:} Covering up is unethical.
	\end{frame}
	
	% =========================================================
	% Q44
	% =========================================================
	\begin{frame}[t]{Q44}
		\small
		\textbf{Question:} Which is the best description of "independence" for internal audit?
		\begin{enumerate}[label=\Alph*.,itemsep=1pt,topsep=2pt]
			\item Freedom from conditions that threaten objective judgment.
			\item The auditor's personal honesty.
			\item The audit committee's meeting schedule.
			\item The audit budget.
		\end{enumerate}
		\vspace{4pt}
		\textbf{Correct Answer: A}
	\end{frame}
	
	\begin{frame}[t]{Q44 -- Explanation}
		\small
		\textbf{Correct: A.} Independence is freedom from conditions that threaten the ability to carry out responsibilities objectively. \\[4pt]
		\textbf{B wrong:} Personal honesty is Integrity. \\[2pt]
		\textbf{C wrong:} Meeting schedules are administrative. \\[2pt]
		\textbf{D wrong:} Budgets are resources.
	\end{frame}
	
	% =========================================================
	% Q45
	% =========================================================
	\begin{frame}[t]{Q45}
		\small
		\textbf{Question:} Which is the best description of "objectivity" for an individual auditor?
		\begin{enumerate}[label=\Alph*.,itemsep=1pt,topsep=2pt]
			\item An unbiased mental attitude.
			\item The audit committee's reporting line.
			\item The audit charter.
			\item The audit plan.
		\end{enumerate}
		\vspace{4pt}
		\textbf{Correct Answer: A}
	\end{frame}
	
	\begin{frame}[t]{Q45 -- Explanation}
		\small
		\textbf{Correct: A.} Objectivity is an individual's unbiased mental attitude. \\[4pt]
		\textbf{B wrong:} Reporting lines are about independence. \\[2pt]
		\textbf{C wrong:} The charter is a document. \\[2pt]
		\textbf{D wrong:} The plan is a planning tool.
	\end{frame}
	
	% =========================================================
	% Q46
	% =========================================================
	\begin{frame}[t]{Q46}
		\small
		\textbf{Question:} An auditor is asked to audit a process she previously operated. She should:
		\begin{enumerate}[label=\Alph*.,itemsep=1pt,topsep=2pt]
			\item Audit it since she knows it best.
			\item Decline and disclose.
			\item Audit only parts she didn't operate.
			\item Delegate without disclosure.
		\end{enumerate}
		\vspace{4pt}
		\textbf{Correct Answer: B}
	\end{frame}
	
	\begin{frame}[t]{Q46 -- Explanation}
		\small
		\textbf{Correct: B.} Operating a process creates a self-review threat; decline and disclose. \\[4pt]
		\textbf{A wrong:} Familiarity creates bias. \\[2pt]
		\textbf{C wrong:} Partial audits still create impairment. \\[2pt]
		\textbf{D wrong:} Delegation without disclosure hides the issue.
	\end{frame}
	
	% =========================================================
	% Q47
	% =========================================================
	\begin{frame}[t]{Q47}
		\small
		\textbf{Question:} Which is NOT a Rule of Conduct under Integrity?
		\begin{enumerate}[label=\Alph*.,itemsep=1pt,topsep=2pt]
			\item Perform work with honesty, diligence, and responsibility.
			\item Observe the law.
			\item Not knowingly be part of illegal activity.
			\item Maintain an unbiased mental attitude.
		\end{enumerate}
		\vspace{4pt}
		\textbf{Correct Answer: D}
	\end{frame}
	
	\begin{frame}[t]{Q47 -- Explanation}
		\small
		\textbf{Correct: D.} Unbiased mental attitude is Objectivity, not Integrity. \\[4pt]
		\textbf{A wrong:} That is Integrity. \\[2pt]
		\textbf{B wrong:} That is Integrity. \\[2pt]
		\textbf{C wrong:} That is Integrity.
	\end{frame}
	
	% =========================================================
	% Q48
	% =========================================================
	\begin{frame}[t]{Q48}
		\small
		\textbf{Question:} Which is NOT a Rule of Conduct under Confidentiality?
		\begin{enumerate}[label=\Alph*.,itemsep=1pt,topsep=2pt]
			\item Be prudent in the use and protection of information.
			\item Not use information for personal benefit.
			\item Not disclose information without authority.
			\item Maintain an unbiased mental attitude.
		\end{enumerate}
		\vspace{4pt}
		\textbf{Correct Answer: D}
	\end{frame}
	
	\begin{frame}[t]{Q48 -- Explanation}
		\small
		\textbf{Correct: D.} Unbiased mental attitude is Objectivity, not Confidentiality. \\[4pt]
		\textbf{A wrong:} That is Confidentiality. \\[2pt]
		\textbf{B wrong:} That is Confidentiality. \\[2pt]
		\textbf{C wrong:} That is Confidentiality.
	\end{frame}
	
	% =========================================================
	% Q49
	% =========================================================
	\begin{frame}[t]{Q49}
		\small
		\textbf{Question:} Which is NOT a Rule of Conduct under Competency?
		\begin{enumerate}[label=\Alph*.,itemsep=1pt,topsep=2pt]
			\item Engage only in services for which you have the necessary knowledge.
			\item Perform services in accordance with the Standards.
			\item Continually improve proficiency.
			\item Not disclose confidential information.
		\end{enumerate}
		\vspace{4pt}
		\textbf{Correct Answer: D}
	\end{frame}
	
	\begin{frame}[t]{Q49 -- Explanation}
		\small
		\textbf{Correct: D.} Not disclosing confidential information is Confidentiality, not Competency. \\[4pt]
		\textbf{A wrong:} That is Competency. \\[2pt]
		\textbf{B wrong:} That is Competency. \\[2pt]
		\textbf{C wrong:} That is Competency.
	\end{frame}
	
	% =========================================================
	% Q50
	% =========================================================
	\begin{frame}[t]{Q50}
		\small
		\textbf{Question:} The IIA Code of Ethics is best described as:
		\begin{enumerate}[label=\Alph*.,itemsep=1pt,topsep=2pt]
			\item A set of principles and rules that guide internal auditors' behavior.
			\item A law with criminal penalties.
			\item A suggestion only.
			\item A marketing tool.
		\end{enumerate}
		\vspace{4pt}
		\textbf{Correct Answer: A}
	\end{frame}
	
	\begin{frame}[t]{Q50 -- Explanation}
		\small
		\textbf{Correct: A.} The Code provides principles and rules of conduct that guide internal auditors' behavior. \\[4pt]
		\textbf{B wrong:} It is not criminal law. \\[2pt]
		\textbf{C wrong:} It is mandatory for those to whom it applies. \\[2pt]
		\textbf{D wrong:} It is not a marketing tool.
	\end{frame}
	
	% =========================================================
	% Q51
	% =========================================================
	\begin{frame}[t]{Q51}
		\small
		\textbf{Question:} An auditor is asked to backdate an audit report. Under Integrity, the auditor should:
		\begin{enumerate}[label=\Alph*.,itemsep=1pt,topsep=2pt]
			\item Backdate it.
			\item Refuse and report.
			\item Backdate it but document the request.
			\item Ask a colleague to do it.
		\end{enumerate}
		\vspace{4pt}
		\textbf{Correct Answer: B}
	\end{frame}
	
	\begin{frame}[t]{Q51 -- Explanation}
		\small
		\textbf{Correct: B.} Backdating is dishonest and possibly illegal; refuse and report. \\[4pt]
		\textbf{A wrong:} Backdating violates Integrity. \\[2pt]
		\textbf{C wrong:} Documenting doesn't cure the violation. \\[2pt]
		\textbf{D wrong:} Delegating wrongdoing is still unethical.
	\end{frame}
	
	% =========================================================
	% Q52
	% =========================================================
	\begin{frame}[t]{Q52}
		\small
		\textbf{Question:} Which is the best description of "material facts" for disclosure under Objectivity?
		\begin{enumerate}[label=\Alph*.,itemsep=1pt,topsep=2pt]
			\item Facts that, if not disclosed, could distort reported information.
			\item All facts.
			\item Only financial facts.
			\item Only facts about family.
		\end{enumerate}
		\vspace{4pt}
		\textbf{Correct Answer: A}
	\end{frame}
	
	\begin{frame}[t]{Q52 -- Explanation}
		\small
		\textbf{Correct: A.} Material facts are those that could distort reporting if withheld. \\[4pt]
		\textbf{B wrong:} Not all facts. \\[2pt]
		\textbf{C wrong:} Non-financial facts can matter. \\[2pt]
		\textbf{D wrong:} Non-family facts can matter.
	\end{frame}
	
	% =========================================================
	% Q53
	% =========================================================
	\begin{frame}[t]{Q53}
		\small
		\textbf{Question:} Which is the best description of a "self-review threat"?
		\begin{enumerate}[label=\Alph*.,itemsep=1pt,topsep=2pt]
			\item The auditor reviews their own prior work or decisions.
			\item The auditor reviews a colleague's work.
			\item The auditor reviews prior-year workpapers.
			\item The auditor uses a standardized program.
		\end{enumerate}
		\vspace{4pt}
		\textbf{Correct Answer: A}
	\end{frame}
	
	\begin{frame}[t]{Q53 -- Explanation}
		\small
		\textbf{Correct: A.} Self-review occurs when an auditor audits a process they designed, implemented, or operated. \\[4pt]
		\textbf{B wrong:} Peer review is a QAIP activity. \\[2pt]
		\textbf{C wrong:} Reviewing prior workpapers provides context. \\[2pt]
		\textbf{D wrong:} Standardized programs are fine.
	\end{frame}
	
	% =========================================================
	% Q54
	% =========================================================
	\begin{frame}[t]{Q54}
		\small
		\textbf{Question:} Which is the best description of a "familiarity threat"?
		\begin{enumerate}[label=\Alph*.,itemsep=1pt,topsep=2pt]
			\item Close relationships that bias judgment.
			\item Financial interest in the client.
			\item Intimidation by management.
			\item Self-review of prior work.
		\end{enumerate}
		\vspace{4pt}
		\textbf{Correct Answer: A}
	\end{frame}
	
	\begin{frame}[t]{Q54 -- Explanation}
		\small
		\textbf{Correct: A.} Familiarity threat arises from close relationships (family, friends) that could bias judgment. \\[4pt]
		\textbf{B wrong:} That's a financial self-interest threat. \\[2pt]
		\textbf{C wrong:} That's an intimidation threat. \\[2pt]
		\textbf{D wrong:} That's a self-review threat.
	\end{frame}
	
	% =========================================================
	% Q55
	% =========================================================
	\begin{frame}[t]{Q55}
		\small
		\textbf{Question:} Which is the best description of an "intimidation threat"?
		\begin{enumerate}[label=\Alph*.,itemsep=1pt,topsep=2pt]
			\item Pressure from management that deters objective action.
			\item Close relationships.
			\item Self-review.
			\item Financial interest.
		\end{enumerate}
		\vspace{4pt}
		\textbf{Correct Answer: A}
	\end{frame}
	
	\begin{frame}[t]{Q55 -- Explanation}
		\small
		\textbf{Correct: A.} Intimidation threat occurs when management pressure deters the auditor from acting objectively. \\[4pt]
		\textbf{B wrong:} That's familiarity. \\[2pt]
		\textbf{C wrong:} That's self-review. \\[2pt]
		\textbf{D wrong:} That's self-interest.
	\end{frame}
	
	% =========================================================
	% Q56
	% =========================================================
	\begin{frame}[t]{Q56}
		\small
		\textbf{Question:} Which is the best description of a "self-interest threat"?
		\begin{enumerate}[label=\Alph*.,itemsep=1pt,topsep=2pt]
			\item Financial or personal interest that biases judgment.
			\item Close relationships.
			\item Self-review.
			\item Intimidation.
		\end{enumerate}
		\vspace{4pt}
		\textbf{Correct Answer: A}
	\end{frame}
	
	\begin{frame}[t]{Q56 -- Explanation}
		\small
		\textbf{Correct: A.} Self-interest threat arises from financial or personal interests that could bias judgment. \\[4pt]
		\textbf{B wrong:} That's familiarity. \\[2pt]
		\textbf{C wrong:} That's self-review. \\[2pt]
		\textbf{D wrong:} That's intimidation.
	\end{frame}
	
	% =========================================================
	% Q57
	% =========================================================
	\begin{frame}[t]{Q57}
		\small
		\textbf{Question:} Which is the best way to manage a threat to objectivity?
		\begin{enumerate}[label=\Alph*.,itemsep=1pt,topsep=2pt]
			\item Ignore it.
			\item Apply safeguards or decline the engagement.
			\item Hide it.
			\item Delegate without disclosure.
		\end{enumerate}
		\vspace{4pt}
		\textbf{Correct Answer: B}
	\end{frame}
	
	\begin{frame}[t]{Q57 -- Explanation}
		\small
		\textbf{Correct: B.} Threats must be eliminated (decline) or reduced to an acceptable level through safeguards (disclosure, reassignment, additional review). \\[4pt]
		\textbf{A wrong:} Ignoring is not acceptable. \\[2pt]
		\textbf{C wrong:} Hiding violates the Code. \\[2pt]
		\textbf{D wrong:} Delegation without disclosure hides the issue.
	\end{frame}
	
	% =========================================================
	% Q58
	% =========================================================
	\begin{frame}[t]{Q58}
		\small
		\textbf{Question:} Which best describes the CAE's role in promoting ethics?
		\begin{enumerate}[label=\Alph*.,itemsep=1pt,topsep=2pt]
			\item Model ethical behavior and support the organization's ethics program.
			\item Write the Code of Ethics.
			\item Enforce criminal law.
			\item Set tax policy.
		\end{enumerate}
		\vspace{4pt}
		\textbf{Correct Answer: A}
	\end{frame}
	
	\begin{frame}[t]{Q58 -- Explanation}
		\small
		\textbf{Correct: A.} The CAE models ethical behavior and supports the organization's ethics program, often through assurance/consulting. \\[4pt]
		\textbf{B wrong:} The IIA writes the Code. \\[2pt]
		\textbf{C wrong:} Criminal enforcement is governmental. \\[2pt]
		\textbf{D wrong:} Tax policy is unrelated.
	\end{frame}
	
	% =========================================================
	% Q59
	% =========================================================
	\begin{frame}[t]{Q59}
		\small
		\textbf{Question:} Which best describes the IIA's guidance on receiving gifts from an audit client?
		\begin{enumerate}[label=\Alph*.,itemsep=1pt,topsep=2pt]
			\item Never accept any gift.
			\item Prohibited if they impair or appear to impair objectivity.
			\item Always accept.
			\item Accept only cash.
		\end{enumerate}
		\vspace{4pt}
		\textbf{Correct Answer: B}
	\end{frame}
	
	\begin{frame}[t]{Q59 -- Explanation}
		\small
		\textbf{Correct: B.} Gifts that impair or appear to impair objectivity are prohibited. \\[4pt]
		\textbf{A wrong:} Not all gifts are prohibited. \\[2pt]
		\textbf{C wrong:} Not always acceptable. \\[2pt]
		\textbf{D wrong:} Cash is especially problematic.
	\end{frame}
	
	% =========================================================
	% Q60
	% =========================================================
	\begin{frame}[t]{Q60}
		\small
		\textbf{Question:} An auditor is asked to audit a process she helped implement six months ago. She should:
		\begin{enumerate}[label=\Alph*.,itemsep=1pt,topsep=2pt]
			\item Proceed.
			\item Decline and disclose.
			\item Audit only the parts she didn't implement.
			\item Delegate without disclosure.
		\end{enumerate}
		\vspace{4pt}
		\textbf{Correct Answer: B}
	\end{frame}
	
	\begin{frame}[t]{Q60 -- Explanation}
		\small
		\textbf{Correct: B.} Implementing a process creates a self-review threat; decline and disclose. \\[4pt]
		\textbf{A wrong:} Six months doesn't remove the impairment. \\[2pt]
		\textbf{C wrong:} Partial audits still create impairment. \\[2pt]
		\textbf{D wrong:} Delegation without disclosure hides the issue.
	\end{frame}
	
	% =========================================================
	% Q61
	% =========================================================
	\begin{frame}[t]{Q61}
		\small
		\textbf{Question:} Which best describes the IIA's position on using confidential information for personal benefit?
		\begin{enumerate}[label=\Alph*.,itemsep=1pt,topsep=2pt]
			\item Permitted if not harmful.
			\item Prohibited.
			\item Permitted if disclosed.
			\item Permitted for family.
		\end{enumerate}
		\vspace{4pt}
		\textbf{Correct Answer: B}
	\end{frame}
	
	\begin{frame}[t]{Q61 -- Explanation}
		\small
		\textbf{Correct: B.} Using confidential information for personal benefit is prohibited under Confidentiality. \\[4pt]
		\textbf{A wrong:} No personal use. \\[2pt]
		\textbf{C wrong:} Disclosure doesn't make it permissible. \\[2pt]
		\textbf{D wrong:} Family use is prohibited.
	\end{frame}
	
	% =========================================================
	% Q62
	% =========================================================
	\begin{frame}[t]{Q62}
		\small
		\textbf{Question:} Which best describes "due care" in the context of ethics?
		\begin{enumerate}[label=\Alph*.,itemsep=1pt,topsep=2pt]
			\item Applying the care expected of a reasonably prudent auditor.
			\item Guaranteeing no errors.
			\item Performing all procedures.
			\item Avoiding all risk.
		\end{enumerate}
		\vspace{4pt}
		\textbf{Correct Answer: A}
	\end{frame}
	
	\begin{frame}[t]{Q62 -- Explanation}
		\small
		\textbf{Correct: A.} Due care is the care and skill expected of a reasonably prudent and competent auditor. \\[4pt]
		\textbf{B wrong:} No guarantee. \\[2pt]
		\textbf{C wrong:} Not all procedures are feasible. \\[2pt]
		\textbf{D wrong:} Risk cannot be eliminated.
	\end{frame}
	
	% =========================================================
	% Q63
	% =========================================================
	\begin{frame}[t]{Q63}
		\small
		\textbf{Question:} Which best describes "professionalism" in the context of internal audit?
		\begin{enumerate}[label=\Alph*.,itemsep=1pt,topsep=2pt]
			\item Adhering to the Code of Ethics and Standards.
			\item Maximizing billable hours.
			\item Avoiding all risk.
			\item Ignoring rules.
		\end{enumerate}
		\vspace{4pt}
		\textbf{Correct Answer: A}
	\end{frame}
	
	\begin{frame}[t]{Q63 -- Explanation}
		\small
		\textbf{Correct: A.} Professionalism means adhering to the Code of Ethics and the Standards, and behaving with integrity, objectivity, confidentiality, and competency. \\[4pt]
		\textbf{B wrong:} Billable hours are irrelevant. \\[2pt]
		\textbf{C wrong:} Risk avoidance is not the goal. \\[2pt]
		\textbf{D wrong:} Ignoring rules is unprofessional.
	\end{frame}
	
	% =========================================================
	% Q64
	% =========================================================
	\begin{frame}[t]{Q64}
		\small
		\textbf{Question:} An auditor is asked to audit a process she previously audited. She should:
		\begin{enumerate}[label=\Alph*.,itemsep=1pt,topsep=2pt]
			\item Proceed; prior audit isn't a conflict.
			\item Decline; prior audit is a conflict.
			\item Audit only new areas.
			\item Delegate.
		\end{enumerate}
		\vspace{4pt}
		\textbf{Correct Answer: A}
	\end{frame}
	
	\begin{frame}[t]{Q64 -- Explanation}
		\small
		\textbf{Correct: A.} Prior audit experience is not a conflict; familiarity from prior audits doesn't impair objectivity. \\[4pt]
		\textbf{B wrong:} Prior audit is not a self-review threat. \\[2pt]
		\textbf{C wrong:} No need to limit scope. \\[2pt]
		\textbf{D wrong:} No need to delegate.
	\end{frame}
	
	% =========================================================
	% Q65
	% =========================================================
	\begin{frame}[t]{Q65}
		\small
		\textbf{Question:} Which best describes "integrity" in practice?
		\begin{enumerate}[label=\Alph*.,itemsep=1pt,topsep=2pt]
			\item Honesty, diligence, responsibility, and law observance.
			\item Bias-free judgment.
			\item Information protection.
			\item Skill maintenance.
		\end{enumerate}
		\vspace{4pt}
		\textbf{Correct Answer: A}
	\end{frame}
	
	\begin{frame}[t]{Q65 -- Explanation}
		\small
		\textbf{Correct: A.} Integrity = honesty, diligence, responsibility, law observance, no knowing involvement in illegal/discreditable activity. \\[4pt]
		\textbf{B wrong:} That's Objectivity. \\[2pt]
		\textbf{C wrong:} That's Confidentiality. \\[2pt]
		\textbf{D wrong:} That's Competency.
	\end{frame}
	
	% =========================================================
	% Q66
	% =========================================================
	\begin{frame}[t]{Q66}
		\small
		\textbf{Question:} An auditor is asked to sign a false statement. Under Integrity, the auditor should:
		\begin{enumerate}[label=\Alph*.,itemsep=1pt,topsep=2pt]
			\item Sign it.
			\item Refuse and report.
			\item Sign but document.
			\item Delegate.
		\end{enumerate}
		\vspace{4pt}
		\textbf{Correct Answer: B}
	\end{frame}
	
	\begin{frame}[t]{Q66 -- Explanation}
		\small
		\textbf{Correct: B.} Signing a false statement violates Integrity; refuse and report. \\[4pt]
		\textbf{A wrong:} Signing is unethical. \\[2pt]
		\textbf{C wrong:} Documenting doesn't cure it. \\[2pt]
		\textbf{D wrong:} Delegating wrongdoing is still unethical.
	\end{frame}
	
	% =========================================================
	% Q67
	% =========================================================
	\begin{frame}[t]{Q67}
		\small
		\textbf{Question:} Which best describes the IIA's position on auditor independence?
		\begin{enumerate}[label=\Alph*.,itemsep=1pt,topsep=2pt]
			\item Independence is required for internal audit activity.
			\item Independence is optional.
			\item Independence is only for external auditors.
			\item Independence is irrelevant.
		\end{enumerate}
		\vspace{4pt}
		\textbf{Correct Answer: A}
	\end{frame}
	
	\begin{frame}[t]{Q7 -- Explanation}
		\small
		\textbf{Correct: A.} Internal audit must be independent, and internal auditors must be objective. \\[4pt]
		\textbf{B wrong:} It is required. \\[2pt]
		\textbf{C wrong:} Internal auditors also need independence/objectivity. \\[2pt]
		\textbf{D wrong:} It is essential.
	\end{frame}
	
	% =========================================================
	% Q68
	% =========================================================
	\begin{frame}[t]{Q68}
		\small
		\textbf{Question:} Which best describes "objectivity" in practice?
		\begin{enumerate}[label=\Alph*.,itemsep=1pt,topsep=2pt]
			\item Unbiased mental attitude, no conflicts, no gifts that impair judgment.
			\item Honesty.
			\item Information protection.
			\item Skill maintenance.
		\end{enumerate}
		\vspace{4pt}
		\textbf{Correct Answer: A}
	\end{frame}
	
	\begin{frame}[t]{Q68 -- Explanation}
		\small
		\textbf{Correct: A.} Objectivity = unbiased mental attitude, no conflicts, no gifts that impair judgment, disclose material facts. \\[4pt]
		\textbf{B wrong:} That's Integrity. \\[2pt]
		\textbf{C wrong:} That's Confidentiality. \\[2pt]
		\textbf{D wrong:} That's Competency.
	\end{frame}
	
	% =========================================================
	% Q69
	% =========================================================
	\begin{frame}[t]{Q69}
		\small
		\textbf{Question:} An auditor is asked to audit a process she previously reviewed as a consultant. She should:
		\begin{enumerate}[label=\Alph*.,itemsep=1pt,topsep=2pt]
			\item Proceed.
			\item Decline and disclose.
			\item Audit only parts she didn't consult on.
			\item Delegate without disclosure.
		\end{enumerate}
		\vspace{4pt}
		\textbf{Correct Answer: B}
	\end{frame}
	
	\begin{frame}[t]{Q69 -- Explanation}
		\small
		\textbf{Correct: B.} Consulting on a process creates a self-review threat; decline and disclose. \\[4pt]
		\textbf{A wrong:} Proceeding violates the Code. \\[2pt]
		\textbf{C wrong:} Partial audits still create impairment. \\[2pt]
		\textbf{D wrong:} Delegation without disclosure hides the issue.
	\end{frame}
	
	% =========================================================
	% Q70
	% =========================================================
	\begin{frame}[t]{Q70}
		\small
		\textbf{Question:} Which best describes "confidentiality" in practice?
		\begin{enumerate}[label=\Alph*.,itemsep=1pt,topsep=2pt]
			\item Prudent use/protection of information; no unauthorized disclosure; no personal benefit.
			\item Honesty.
			\item Unbiased judgment.
			\item Skill maintenance.
		\end{enumerate}
		\vspace{4pt}
		\textbf{Correct Answer: A}
	\end{frame}
	
	\begin{frame}[t]{Q70 -- Explanation}
		\small
		\textbf{Correct: A.} Confidentiality = prudent use/protection, no unauthorized disclosure, no personal benefit. \\[4pt]
		\textbf{B wrong:} That's Integrity. \\[2pt]
		\textbf{C wrong:} That's Objectivity. \\[2pt]
		\textbf{D wrong:} That's Competency.
	\end{frame}
	
	% =========================================================
	% Q71
	% =========================================================
	\begin{frame}[t]{Q71}
		\small
		\textbf{Question:} Which best describes "competency" in practice?
		\begin{enumerate}[label=\Alph*.,itemsep=1pt,topsep=2pt]
			\item Knowledge, skills, experience; CPD; engage only in services you can perform.
			\item Honesty.
			\item Unbiased judgment.
			\item Information protection.
		\end{enumerate}
		\vspace{4pt}
		\textbf{Correct Answer: A}
	\end{frame}
	
	\begin{frame}[t]{Q71 -- Explanation}
		\small
		\textbf{Correct: A.} Competency = knowledge, skills, experience, CPD, and only engaging in services you can perform. \\[4pt]
		\textbf{B wrong:} That's Integrity. \\[2pt]
		\textbf{C wrong:} That's Objectivity. \\[2pt]
		\textbf{D wrong:} That's Confidentiality.
	\end{frame}
	
	% =========================================================
	% Q72
	% =========================================================
	\begin{frame}[t]{Q72}
		\small
		\textbf{Question:} An auditor is asked to audit a process she previously managed but no longer manages. She should:
		\begin{enumerate}[label=\Alph*.,itemsep=1pt,topsep=2pt]
			\item Proceed.
			\item Decline and disclose.
			\item Audit only new areas.
			\item Delegate without disclosure.
		\end{enumerate}
		\vspace{4pt}
		\textbf{Correct Answer: B}
	\end{frame}
	
	\begin{frame}[t]{Q72 -- Explanation}
		\small
		\textbf{Correct: B.} Prior management of a process creates a self-review threat; decline and disclose. \\[4pt]
		\textbf{A wrong:} Proceeding violates the Code. \\[2pt]
		\textbf{C wrong:} Partial audits still create impairment. \\[2pt]
		\textbf{D wrong:} Delegation without disclosure hides the issue.
	\end{frame}
	
	% =========================================================
	% Q73
	% =========================================================
	\begin{frame}[t]{Q73}
		\small
		\textbf{Question:} Which best describes the IIA's position on auditor objectivity?
		\begin{enumerate}[label=\Alph*.,itemsep=1pt,topsep=2pt]
			\item Objectivity is required of all internal auditors.
			\item Objectivity is optional.
			\item Objectivity is only for CAEs.
			\item Objectivity is irrelevant.
		\end{enumerate}
		\vspace{4pt}
		\textbf{Correct Answer: A}
	\end{frame}
	
	\begin{frame}[t]{Q73 -- Explanation}
		\small
		\textbf{Correct: A.} All internal auditors must be objective. \\[4pt]
		\textbf{B wrong:} It is required. \\[2pt]
		\textbf{C wrong:} It applies to all internal auditors. \\[2pt]
		\textbf{D wrong:} It is essential.
	\end{frame}
	
	% =========================================================
	% Q74
	% =========================================================
	\begin{frame}[t]{Q74}
		\small
		\textbf{Question:} An auditor is asked to audit a process she previously designed. She should:
		\begin{enumerate}[label=\Alph*.,itemsep=1pt,topsep=2pt]
			\item Proceed.
			\item Decline and disclose.
			\item Audit only parts she didn't design.
			\item Delegate without disclosure.
		\end{enumerate}
		\vspace{4pt}
		\textbf{Correct Answer: B}
	\end{frame}
	
	\begin{frame}[t]{Q74 -- Explanation}
		\small
		\textbf{Correct: B.} Designing a process creates a self-review threat; decline and disclose. \\[4pt]
		\textbf{A wrong:} Proceeding violates the Code. \\[2pt]
		\textbf{C wrong:} Partial audits still create impairment. \\[2pt]
		\textbf{D wrong:} Delegation without disclosure hides the issue.
	\end{frame}
	
	% =========================================================
	% Q75
	% =========================================================
	\begin{frame}[t]{Q75}
		\small
		\textbf{Question:} Which best describes the IIA's position on auditor confidentiality?
		\begin{enumerate}[label=\Alph*.,itemsep=1pt,topsep=2pt]
			\item Confidentiality is required of all internal auditors.
			\item Confidentiality is optional.
			\item Confidentiality is only for CAEs.
			\item Confidentiality is irrelevant.
		\end{enumerate}
		\vspace{4pt}
		\textbf{Correct Answer: A}
	\end{frame}
	
	\begin{frame}[t]{Q75 -- Explanation}
		\small
		\textbf{Correct: A.} All internal auditors must maintain confidentiality. \\[4pt]
		\textbf{B wrong:} It is required. \\[2pt]
		\textbf{C wrong:} It applies to all internal auditors. \\[2pt]
		\textbf{D wrong:} It is essential.
	\end{frame}
	
	% =========================================================
	% Q76
	% =========================================================
	\begin{frame}[t]{Q76}
		\small
		\textbf{Question:} An auditor is asked to audit a process she previously operated. She should:
		\begin{enumerate}[label=\Alph*.,itemsep=1pt,topsep=2pt]
			\item Proceed.
			\item Decline and disclose.
			\item Audit only new areas.
			\item Delegate without disclosure.
		\end{enumerate}
		\vspace{4pt}
		\textbf{Correct Answer: B}
	\end{frame}
	
	\begin{frame}[t]{Q76 -- Explanation}
		\small
		\textbf{Correct: B.} Operating a process creates a self-review threat; decline and disclose. \\[4pt]
		\textbf{A wrong:} Proceeding violates the Code. \\[2pt]
		\textbf{C wrong:} Partial audits still create impairment. \\[2pt]
		\textbf{D wrong:} Delegation without disclosure hides the issue.
	\end{frame}
	
	% =========================================================
	% Q77
	% =========================================================
	\begin{frame}[t]{Q77}
		\small
		\textbf{Question:} Which best describes the IIA's position on auditor competency?
		\begin{enumerate}[label=\Alph*.,itemsep=1pt,topsep=2pt]
			\item Competency is required of all internal auditors.
			\item Competency is optional.
			\item Competency is only for CAEs.
			\item Competency is irrelevant.
		\end{enumerate}
		\vspace{4pt}
		\textbf{Correct Answer: A}
	\end{frame}
	
	\begin{frame}[t]{Q77 -- Explanation}
		\small
		\textbf{Correct: A.} All internal auditors must be competent. \\[4pt]
		\textbf{B wrong:} It is required. \\[2pt]
		\textbf{C wrong:} It applies to all internal auditors. \\[2pt]
		\textbf{D wrong:} It is essential.
	\end{frame}
	
	% =========================================================
	% Q78
	% =========================================================
	\begin{frame}[t]{Q78}
		\small
		\textbf{Question:} An auditor is asked to audit a process she previously implemented. She should:
		\begin{enumerate}[label=\Alph*.,itemsep=1pt,topsep=2pt]
			\item Proceed.
			\item Decline and disclose.
			\item Audit only new areas.
			\item Delegate without disclosure.
		\end{enumerate}
		\vspace{4pt}
		\textbf{Correct Answer: B}
	\end{frame}
	
	\begin{frame}[t]{Q78 -- Explanation}
		\small
		\textbf{Correct: B.} Implementing a process creates a self-review threat; decline and disclose. \\[4pt]
		\textbf{A wrong:} Proceeding violates the Code. \\[2pt]
		\textbf{C wrong:} Partial audits still create impairment. \\[2pt]
		\textbf{D wrong:} Delegation without disclosure hides the issue.
	\end{frame}
	
	% =========================================================
	% Q79
	% =========================================================
	\begin{frame}[t]{Q79}
		\small
		\textbf{Question:} Which best describes the IIA's position on auditor integrity?
		\begin{enumerate}[label=\Alph*.,itemsep=1pt,topsep=2pt]
			\item Integrity is required of all internal auditors.
			\item Integrity is optional.
			\item Integrity is only for CAEs.
			\item Integrity is irrelevant.
		\end{enumerate}
		\vspace{4pt}
		\textbf{Correct Answer: A}
	\end{frame}
	
	\begin{frame}[t]{Q79 -- Explanation}
		\small
		\textbf{Correct: A.} All internal auditors must act with integrity. \\[4pt]
		\textbf{B wrong:} It is required. \\[2pt]
		\textbf{C wrong:} It applies to all internal auditors. \\[2pt]
		\textbf{D wrong:} It is essential.
	\end{frame}
	
	% =========================================================
	% Q80
	% =========================================================
	\begin{frame}[t]{Q80}
		\small
		\textbf{Question:} An auditor is asked to audit a process she previously reviewed as a peer reviewer. She should:
		\begin{enumerate}[label=\Alph*.,itemsep=1pt,topsep=2pt]
			\item Proceed.
			\item Decline and disclose.
			\item Audit only new areas.
			\item Delegate without disclosure.
		\end{enumerate}
		\vspace{4pt}
		\textbf{Correct Answer: A}
	\end{frame}
	
	\begin{frame}[t]{Q80 -- Explanation}
		\small
		\textbf{Correct: A.} Peer review is a QAIP activity and does not impair objectivity. \\[4pt]
		\textbf{B wrong:} Peer review is not a self-review threat. \\[2pt]
		\textbf{C wrong:} No need to limit scope. \\[2pt]
		\textbf{D wrong:} No need to delegate.
	\end{frame}
	
	% =========================================================
	% Q81
	% =========================================================
	\begin{frame}[t]{Q81}
		\small
		\textbf{Question:} Which best describes the IIA's position on auditor CPD?
		\begin{enumerate}[label=\Alph*.,itemsep=1pt,topsep=2pt]
			\item CPD is required of all internal auditors.
			\item CPD is optional.
			\item CPD is only for CAEs.
			\item CPD is irrelevant.
		\end{enumerate}
		\vspace{4pt}
		\textbf{Correct Answer: A}
	\end{frame}
	
	\begin{frame}[t]{Q81 -- Explanation}
		\small
		\textbf{Correct: A.} All internal auditors must maintain and enhance competence through CPD. \\[4pt]
		\textbf{B wrong:} It is required. \\[2pt]
		\textbf{C wrong:} It applies to all internal auditors. \\[2pt]
		\textbf{D wrong:} It is essential.
	\end{frame}
	
	% =========================================================
	% Q82
	% =========================================================
	\begin{frame}[t]{Q82}
		\small
		\textbf{Question:} An auditor is asked to audit a process she previously approved as a manager. She should:
		\begin{enumerate}[label=\Alph*.,itemsep=1pt,topsep=2pt]
			\item Proceed.
			\item Decline and disclose.
			\item Audit only new areas.
			\item Delegate without disclosure.
		\end{enumerate}
		\vspace{4pt}
		\textbf{Correct Answer: B}
	\end{frame}
	
	\begin{frame}[t]{Q82 -- Explanation}
		\small
		\textbf{Correct: B.} Approving a process as manager creates a self-review threat; decline and disclose. \\[4pt]
		\textbf{A wrong:} Proceeding violates the Code. \\[2pt]
		\textbf{C wrong:} Partial audits still create impairment. \\[2pt]
		\textbf{D wrong:} Delegation without disclosure hides the issue.
	\end{frame}
	
	% =========================================================
	% Q83
	% =========================================================
	\begin{frame}[t]{Q83}
		\small
		\textbf{Question:} Which best describes the IIA's position on auditor skepticism?
		\begin{enumerate}[label=\Alph*.,itemsep=1pt,topsep=2pt]
			\item A questioning mind and critical assessment of evidence.
			\item Assuming management is dishonest.
			\item Ignoring evidence.
			\item Relying on prior conclusions.
		\end{enumerate}
		\vspace{4pt}
		\textbf{Correct Answer: A}
	\end{frame}
	
	\begin{frame}[t]{Q83 -- Explanation}
		\small
		\textbf{Correct: A.} Professional skepticism is a questioning mind, alertness to conditions suggesting misstatement, and critical evaluation of evidence. \\[4pt]
		\textbf{B wrong:} It is not cynicism. \\[2pt]
		\textbf{C wrong:} Ignoring evidence is the opposite. \\[2pt]
		\textbf{D wrong:} Prior conclusions must be re-validated.
	\end{frame}
	
	% =========================================================
	% Q84
	% =========================================================
	\begin{frame}[t]{Q84}
		\small
		\textbf{Question:} An auditor is asked to audit a process she previously tested as an internal auditor. She should:
		\begin{enumerate}[label=\Alph*.,itemsep=1pt,topsep=2pt]
			\item Proceed.
			\item Decline and disclose.
			\item Audit only new areas.
			\item Delegate without disclosure.
		\end{enumerate}
		\vspace{4pt}
		\textbf{Correct Answer: A}
	\end{frame}
	
	\begin{frame}[t]{Q84 -- Explanation}
		\small
		\textbf{Correct: A.} Prior testing as an internal auditor is not a self-review threat. \\[4pt]
		\textbf{B wrong:} No conflict. \\[2pt]
		\textbf{C wrong:} No need to limit scope. \\[2pt]
		\textbf{D wrong:} No need to delegate.
	\end{frame}
	
	% =========================================================
	% Q85
	% =========================================================
	\begin{frame}[t]{Q85}
		\small
		\textbf{Question:} Which best describes the IIA's position on auditor objectivity when auditing a friend?
		\begin{enumerate}[label=\Alph*.,itemsep=1pt,topsep=2pt]
			\item Objectivity may be impaired; disclose and consider reassignment.
			\item Objectivity is not impaired.
			\item Objectivity is always impaired.
			\item Objectivity is irrelevant.
		\end{enumerate}
		\vspace{4pt}
		\textbf{Correct Answer: A}
	\end{frame}
	
	\begin{frame}[t]{Q85 -- Explanation}
		\small
		\textbf{Correct: A.} Close relationships may impair objectivity; disclose and consider reassignment. \\[4pt]
		\textbf{B wrong:} Friendship can bias. \\[2pt]
		\textbf{C wrong:} Not always impaired, but a threat exists. \\[2pt]
		\textbf{D wrong:} Objectivity is always relevant.
	\end{frame}
	
	% =========================================================
	% Q86
	% =========================================================
	\begin{frame}[t]{Q86}
		\small
		\textbf{Question:} An auditor is asked to audit a process she previously reviewed as an external consultant. She should:
		\begin{enumerate}[label=\Alph*.,itemsep=1pt,topsep=2pt]
			\item Proceed.
			\item Decline and disclose.
			\item Audit only new areas.
			\item Delegate without disclosure.
		\end{enumerate}
		\vspace{4pt}
		\textbf{Correct Answer: B}
	\end{frame}
	
	\begin{frame}[t]{Q86 -- Explanation}
		\small
		\textbf{Correct: B.} Consulting creates a self-review threat; decline and disclose. \\[4pt]
		\textbf{A wrong:} Proceeding violates the Code. \\[2pt]
		\textbf{C wrong:} Partial audits still create impairment. \\[2pt]
		\textbf{D wrong:} Delegation without disclosure hides the issue.
	\end{frame}
	
	% =========================================================
	% Q87
	% =========================================================
	\begin{frame}[t]{Q87}
		\small
		\textbf{Question:} Which best describes the IIA's position on auditor objectivity when auditing a family member?
		\begin{enumerate}[label=\Alph*.,itemsep=1pt,topsep=2pt]
			\item Objectivity may be impaired; disclose and decline.
			\item Objectivity is not impaired.
			\item Objectivity is always impaired.
			\item Objectivity is irrelevant.
		\end{enumerate}
		\vspace{4pt}
		\textbf{Correct Answer: A}
	\end{frame}
	
	\begin{frame}[t]{Q87 -- Explanation}
		\small
		\textbf{Correct: A.} Family relationships impair objectivity; disclose and decline. \\[4pt]
		\textbf{B wrong:} Family can bias. \\[2pt]
		\textbf{C wrong:} Impairment is likely but the standard is disclosure. \\[2pt]
		\textbf{D wrong:} Objectivity is always relevant.
	\end{frame}
	
	% =========================================================
	% Q88
	% =========================================================
	\begin{frame}[t]{Q88}
		\small
		\textbf{Question:} An auditor is asked to audit a process she previously reviewed as a quality assurance reviewer. She should:
		\begin{enumerate}[label=\Alph*.,itemsep=1pt,topsep=2pt]
			\item Proceed.
			\item Decline and disclose.
			\item Audit only new areas.
			\item Delegate without disclosure.
		\end{enumerate}
		\vspace{4pt}
		\textbf{Correct Answer: A}
	\end{frame}
	
	\begin{frame}[t]{Q88 -- Explanation}
		\small
		\textbf{Correct: A.} QA review is a QAIP activity and does not impair objectivity. \\[4pt]
		\textbf{B wrong:} No conflict. \\[2pt]
		\textbf{C wrong:} No need to limit scope. \\[2pt]
		\textbf{D wrong:} No need to delegate.
	\end{frame}
	
	% =========================================================
	% Q89
	% =========================================================
	\begin{frame}[t]{Q89}
		\small
		\textbf{Question:} Which best describes the IIA's position on auditor objectivity when auditing a process the auditor previously approved?
		\begin{enumerate}[label=\Alph*.,itemsep=1pt,topsep=2pt]
			\item Objectivity may be impaired; disclose and decline.
			\item Objectivity is not impaired.
			\item Objectivity is always impaired.
			\item Objectivity is irrelevant.
		\end{enumerate}
		\vspace{4pt}
		\textbf{Correct Answer: A}
	\end{frame}
	
	\begin{frame}[t]{Q89 -- Explanation}
		\small
		\textbf{Correct: A.} Approving a process creates a self-review threat; disclose and decline. \\[4pt]
		\textbf{B wrong:} Approval creates bias. \\[2pt]
		\textbf{C wrong:} Impairment is likely but the standard is disclosure. \\[2pt]
		\textbf{D wrong:} Objectivity is always relevant.
	\end{frame}
	
	% =========================================================
	% Q90
	% =========================================================
	\begin{frame}[t]{Q90}
		\small
		\textbf{Question:} An auditor is asked to audit a process she previously reviewed as a fraud investigator. She should:
		\begin{enumerate}[label=\Alph*.,itemsep=1pt,topsep=2pt]
			\item Proceed.
			\item Decline and disclose.
			\item Audit only new areas.
			\item Delegate without disclosure.
		\end{enumerate}
		\vspace{4pt}
		\textbf{Correct Answer: A}
	\end{frame}
	
	\begin{frame}[t]{Q90 -- Explanation}
		\small
		\textbf{Correct: A.} Prior fraud investigation is not a self-review threat; the auditor can proceed. \\[4pt]
		\textbf{B wrong:} No conflict. \\[2pt]
		\textbf{C wrong:} No need to limit scope. \\[2pt]
		\textbf{D wrong:} No need to delegate.
	\end{frame}
	
	% =========================================================
	% Q91
	% =========================================================
	\begin{frame}[t]{Q91}
		\small
		\textbf{Question:} Which best describes the IIA's position on auditor objectivity when auditing a process the auditor previously designed?
		\begin{enumerate}[label=\Alph*.,itemsep=1pt,topsep=2pt]
			\item Objectivity may be impaired; disclose and decline.
			\item Objectivity is not impaired.
			\item Objectivity is always impaired.
			\item Objectivity is irrelevant.
		\end{enumerate}
		\vspace{4pt}
		\textbf{Correct Answer: A}
	\end{frame}
	
	\begin{frame}[t]{Q91 -- Explanation}
		\small
		\textbf{Correct: A.} Designing creates a self-review threat; disclose and decline. \\[4pt]
		\textbf{B wrong:} Design creates bias. \\[2pt]
		\textbf{C wrong:} Impairment is likely but the standard is disclosure. \\[2pt]
		\textbf{D wrong:} Objectivity is always relevant.
	\end{frame}
	
	% =========================================================
	% Q92
	% =========================================================
	\begin{frame}[t]{Q92}
		\small
		\textbf{Question:} An auditor is asked to audit a process she previously reviewed as a peer. She should:
		\begin{enumerate}[label=\Alph*.,itemsep=1pt,topsep=2pt]
			\item Proceed.
			\item Decline and disclose.
			\item Audit only new areas.
			\item Delegate without disclosure.
		\end{enumerate}
		\vspace{4pt}
		\textbf{Correct Answer: A}
	\end{frame}
	
	\begin{frame}[t]{Q92 -- Explanation}
		\small
		\textbf{Correct: A.} Peer review is a QAIP activity and does not impair objectivity. \\[4pt]
		\textbf{B wrong:} No conflict. \\[2pt]
		\textbf{C wrong:} No need to limit scope. \\[2pt]
		\textbf{D wrong:} No need to delegate.
	\end{frame}
	
	% =========================================================
	% Q93
	% =========================================================
	\begin{frame}[t]{Q93}
		\small
		\textbf{Question:} Which best describes the IIA's position on auditor objectivity when auditing a process the auditor previously implemented?
		\begin{enumerate}[label=\Alph*.,itemsep=1pt,topsep=2pt]
			\item Objectivity may be impaired; disclose and decline.
			\item Objectivity is not impaired.
			\item Objectivity is always impaired.
			\item Objectivity is irrelevant.
		\end{enumerate}
		\vspace{4pt}
		\textbf{Correct Answer: A}
	\end{frame}
	
	\begin{frame}[t]{Q93 -- Explanation}
		\small
		\textbf{Correct: A.} Implementing creates a self-review threat; disclose and decline. \\[4pt]
		\textbf{B wrong:} Implementation creates bias. \\[2pt]
		\textbf{C wrong:} Impairment is likely but the standard is disclosure. \\[2pt]
		\textbf{D wrong:} Objectivity is always relevant.
	\end{frame}
	
	% =========================================================
	% Q94
	% =========================================================
	\begin{frame}[t]{Q94}
		\small
		\textbf{Question:} An auditor is asked to audit a process she previously reviewed as an external auditor. She should:
		\begin{enumerate}[label=\Alph*.,itemsep=1pt,topsep=2pt]
			\item Proceed.
			\item Decline and disclose.
			\item Audit only new areas.
			\item Delegate without disclosure.
		\end{enumerate}
		\vspace{4pt}
		\textbf{Correct Answer: A}
	\end{frame}
	
	\begin{frame}[t]{Q94 -- Explanation}
		\small
		\textbf{Correct: A.} Prior external audit is not a self-review threat for internal audit. \\[4pt]
		\textbf{B wrong:} No conflict. \\[2pt]
		\textbf{C wrong:} No need to limit scope. \\[2pt]
		\textbf{D wrong:} No need to delegate.
	\end{frame}
	
	% =========================================================
	% Q95
	% =========================================================
	\begin{frame}[t]{Q95}
		\small
		\textbf{Question:} Which best describes the IIA's position on auditor objectivity when auditing a process the auditor previously operated?
		\begin{enumerate}[label=\Alph*.,itemsep=1pt,topsep=2pt]
			\item Objectivity may be impaired; disclose and decline.
			\item Objectivity is not impaired.
			\item Objectivity is always impaired.
			\item Objectivity is irrelevant.
		\end{enumerate}
		\vspace{4pt}
		\textbf{Correct Answer: A}
	\end{frame}
	
	\begin{frame}[t]{Q95 -- Explanation}
		\small
		\textbf{Correct: A.} Operating creates a self-review threat; disclose and decline. \\[4pt]
		\textbf{B wrong:} Operation creates bias. \\[2pt]
		\textbf{C wrong:} Impairment is likely but the standard is disclosure. \\[2pt]
		\textbf{D wrong:} Objectivity is always relevant.
	\end{frame}
	
	% =========================================================
	% Q96
	% =========================================================
	\begin{frame}[t]{Q96}
		\small
		\textbf{Question:} An auditor is asked to audit a process she previously reviewed as a compliance officer. She should:
		\begin{enumerate}[label=\Alph*.,itemsep=1pt,topsep=2pt]
			\item Proceed.
			\item Decline and disclose.
			\item Audit only new areas.
			\item Delegate without disclosure.
		\end{enumerate}
		\vspace{4pt}
		\textbf{Correct Answer: B}
	\end{frame}
	
	\begin{frame}[t]{Q96 -- Explanation}
		\small
		\textbf{Correct: B.} Compliance monitoring creates a self-review threat; decline and disclose. \\[4pt]
		\textbf{A wrong:} Proceeding violates the Code. \\[2pt]
		\textbf{C wrong:} Partial audits still create impairment. \\[2pt]
		\textbf{D wrong:} Delegation without disclosure hides the issue.
	\end{frame}
	
	% =========================================================
	% Q97
	% =========================================================
	\begin{frame}[t]{Q97}
		\small
		\textbf{Question:} Which best describes the IIA's position on auditor objectivity when auditing a process the auditor previously managed?
		\begin{enumerate}[label=\Alph*.,itemsep=1pt,topsep=2pt]
			\item Objectivity may be impaired; disclose and decline.
			\item Objectivity is not impaired.
			\item Objectivity is always impaired.
			\item Objectivity is irrelevant.
		\end{enumerate}
		\vspace{4pt}
		\textbf{Correct Answer: A}
	\end{frame}
	
	\begin{frame}[t]{Q97 -- Explanation}
		\small
		\textbf{Correct: A.} Managing creates a self-review threat; disclose and decline. \\[4pt]
		\textbf{B wrong:} Management creates bias. \\[2pt]
		\textbf{C wrong:} Impairment is likely but the standard is disclosure. \\[2pt]
		\textbf{D wrong:} Objectivity is always relevant.
	\end{frame}
	
	% =========================================================
	% Q98
	% =========================================================
	\begin{frame}[t]{Q98}
		\small
		\textbf{Question:} An auditor is asked to audit a process she previously reviewed as a peer reviewer for the IIA. She should:
		\begin{enumerate}[label=\Alph*.,itemsep=1pt,topsep=2pt]
			\item Proceed.
			\item Decline and disclose.
			\item Audit only new areas.
			\item Delegate without disclosure.
		\end{enumerate}
		\vspace{4pt}
		\textbf{Correct Answer: A}
	\end{frame}
	
	\begin{frame}[t]{Q98 -- Explanation}
		\small
		\textbf{Correct: A.} IIA peer review is a QAIP activity and does not impair objectivity. \\[4pt]
		\textbf{B wrong:} No conflict. \\[2pt]
		\textbf{C wrong:} No need to limit scope. \\[2pt]
		\textbf{D wrong:} No need to delegate.
	\end{frame}
	
	% =========================================================
	% Q99
	% =========================================================
	\begin{frame}[t]{Q99}
		\small
		\textbf{Question:} Which best describes the IIA's position on auditor objectivity when auditing a process the auditor previously reviewed as an internal quality reviewer?
		\begin{enumerate}[label=\Alph*.,itemsep=1pt,topsep=2pt]
			\item Objectivity may be impaired; disclose and decline.
			\item Objectivity is not impaired.
			\item Objectivity is always impaired.
			\item Objectivity is irrelevant.
		\end{enumerate}
		\vspace{4pt}
		\textbf{Correct Answer: B}
	\end{frame}
	
	\begin{frame}[t]{Q99 -- Explanation}
		\small
		\textbf{Correct: B.} Internal QA review is a QAIP activity and does not impair objectivity. \\[4pt]
		\textbf{A wrong:} QA review is not a self-review threat. \\[2pt]
		\textbf{C wrong:} Not always impaired. \\[2pt]
		\textbf{D wrong:} Objectivity is always relevant.
	\end{frame}
	
	% =========================================================
	% Q100
	% =========================================================
	\begin{frame}[t]{Q100}
		\small
		\textbf{Question:} Which of the following best summarizes the IIA Code of Ethics?
		\begin{enumerate}[label=\Alph*.,itemsep=1pt,topsep=2pt]
			\item Four principles (Integrity, Objectivity, Confidentiality, Competency) with Rules of Conduct.
			\item Five principles with no rules.
			\item Three principles with laws.
			\item Two principles with regulations.
		\end{enumerate}
		\vspace{4pt}
		\textbf{Correct Answer: A}
	\end{frame}
	
	\begin{frame}[t]{Q100 -- Explanation}
		\small
		\textbf{Correct: A.} The Code has four principles and corresponding Rules of Conduct. \\[4pt]
		\textbf{B wrong:} Four principles, not five. \\[2pt]
		\textbf{C wrong:} Four, not three. \\[2pt]
		\textbf{D wrong:} Four, not two.
	\end{frame}
	
	\begin{frame}{End}
		\centering
		\Large End of 100-Question Ethics \& Professionalism Deck\\[8pt]
		\normalsize Review weak areas and re-test.\\[4pt]
		\small Add more questions by copying the Q/Explanation frame pairs.
	\end{frame}
	
\end{document}